\chapter{Виходи машини: як мозок керує м'язами}
\chaptermark{Виходи машини}
\label{sec:muscles}

\section{Швидкий вихід}

Усе, що машина робить у світі швидко, вона робить м'язами. Слово, погляд, крок, усмішка --- це скорочення м'язів. Другий, повільний вихід --- хімію тіла --- розібрано в розділі~\ref{sec:chemoutput}. На рис.~\ref{fig:machine} вихід був стрілкою «м'язи»; тепер подивімося, що за цією стрілкою.

Остання ланка кола --- \nt{ACET}-нейрони, ті самі, в яких у таблиці~\ref{tab:types} записано «вихід на м'яз». Кожне м'язове волокно отримує вхід рівно від одного такого нейрона, а один нейрон керує групою волокон --- від кількох сотень до пари тисяч~\cite{Feinstein1955mu}. У м'язах, яким потрібне тонке підлаштування, одиниці дрібніші, у великих м'язах ніг --- більші. Нейрон разом зі своїми волокнами називають \emph{руховою одиницею}: це найменший шматок м'яза, який мозок може ввімкнути окремо.

Синапс нейрона на м'язі влаштований зовсім не так, як синапси кори з розділу~\ref{sec:code}. Там більшість імпульсів губилася. Тут кожен імпульс вивільняє в кілька разів більше ацетилхоліну, ніж потрібно, щоб волокно спрацювало~\cite{WoodSlater2001}. Це синапс із \emph{запасом надійності}: один імпульс нейрона --- рівно одне скорочення волокон. Усередині машини можна дозволити собі ненадійні зв'язки й виправляти помилки надлишковістю; на виході помилка коштує руху, і надійність закладено просто в залізо.

\section{Сила: частота й кількість одиниць}

Один імпульс дає поодиноке скорочення --- посмикування: сила наростає за кілька десятків мілісекунд і спадає. Добре наближення для нього~\cite{Fuglevand1993}
\begin{equation}
h(t) \;=\; F_1\,\frac{t}{T_c}\,e^{\,1-t/T_c},
\label{eq:twitch}
\end{equation}
де $T_c$ --- час наростання, порядку $50\ms$, а $F_1$ --- сила одного посмикування. Якщо імпульси йдуть частіше, посмикування додаються:
\begin{empheq}[box=\keybox]{equation}
F(t) \;=\; \sum_k h\bigl(t-t_k\bigr) .
\label{eq:forcesum}
\end{empheq}
Це той самий фільтр нижніх частот, що перетворював імпульси на концентрацію в розділі~\ref{sec:serocomputer}, тільки тепер на виході не концентрація, а сила. М'яз --- цифро-аналоговий перетворювач з імпульсів у силу (рис.~\ref{fig:tetanus}). У нашій моделі за п'яти імпульсів на секунду посмикування майже не перекриваються й сила стрибає від $0.2F_1$ до $1.1F_1$; за п'ятнадцяти вона вже тримається між $1.8F_1$ і $2.2F_1$; за сорока стає майже рівною, близько $5.4F_1$. Справжній м'яз за частих імпульсів насичується, тож лінійна сума~\eqref{eq:forcesum} правильна лише до кількох десятків імпульсів за секунду.

\begin{figure}[t]
\centering
\begin{tikzpicture}[>=Stealth,x=20cm,y=0.55cm]
\draw[->,gray!70!black] (0,0) -- (0.66,0) node[right,font=\small,black] {$t$, с};
\draw[->,gray!70!black] (0,0) -- (0,6.4) node[left,font=\small,black] {$F/F_1$};
\foreach \t in {0,0.2,0.4,0.6}{\draw[gray!60!black] (\t,-0.1) -- (\t,0.1); \node[below,font=\scriptsize] at (\t,0) {\t};}
\foreach \f in {1,3,5}{\draw[gray!60!black] (-0.004,\f) -- (0.004,\f); \node[left,font=\scriptsize] at (0,\f) {\f};}
\draw[thick,blue!60!black] plot file {figures/muscle_twitch_5.dat};
\draw[thick,green!45!black] plot file {figures/muscle_twitch_15.dat};
\draw[thick,red!70!black] plot file {figures/muscle_twitch_40.dat};
\node[font=\scriptsize,blue!60!black,anchor=west] at (0.605,0.9) {5 імп/с};
\node[font=\scriptsize,green!45!black,anchor=west] at (0.605,2.1) {15 імп/с};
\node[font=\scriptsize,red!70!black,anchor=west] at (0.605,5.4) {40 імп/с};
\end{tikzpicture}
\caption{М'яз як цифро-аналоговий перетворювач: сила~\eqref{eq:forcesum} за регулярних імпульсів однієї рухової одиниці, $T_c=50\ms$. Рідкі імпульси дають окремі посмикування, часті зливаються в рівну силу --- так само, як часті імпульси зливаються в рівну концентрацію на рис.~\ref{fig:lowpass}.}
\label{fig:tetanus}
\end{figure}

У мозку дві ручки сили: частота імпульсів кожної одиниці й кількість увімкнених одиниць. І друга ручка влаштована дуже винахідливо. Одиниці в м'язі різні: найслабші в сотню разів слабші за найсильніші. Коли потрібна дедалі більша сила, одиниці вмикаються \emph{за порядком розміру}, від малих до великих~\cite{Henneman1957}. Ніхто цього порядку не обчислює: малі нейрони просто легше збуджуються від того самого спільного входу, тож порядок увімкнення задано самим залізом.

Що це дає? Нехай кожна наступна за порядком одиниця сильніша за попередню в $q$ разів, де $q$ трохи більше за одиницю. Сила всіх увімкнених одиниць --- сума геометричної прогресії, і майже вся вона припадає на останні одиниці. Тоді чергова ввімкнена одиниця додає
\begin{empheq}[box=\keybox]{equation}
\frac{\Delta F}{F} \;\approx\; q-1 \;=\; \text{const} ,
\label{eq:sizeprinciple}
\end{empheq}
тобто крок сили пропорційний самій силі. За слабкого зусилля мозок дозує його дрібними порціями, за сильного --- великими. Це \emph{число з рухомою комою}: порядок задає кількість увімкнених одиниць, а мантису --- частота їхніх імпульсів. Та сама логарифмічна шкала, що в датчиків розділу~\ref{sec:sensors}: вхід і вихід машини міряють світ у відносних частках.

Рухома кома має зворотний бік. Розкид сили теж зростає пропорційно їй самій: сильний рух неминуче менш точний, ніж слабкий. За однією впливовою гіпотезою, саме цей шум, що залежить від сигналу, пояснює, чому наші рухи плавні: плавна команда дає найменший розкид у кінцевій точці~\cite{HarrisWolpert1998}.

\section{Тягнути, але не штовхати: два проводи на суглоб}

М'яз уміє лише тягнути. Штовхати він не може, як не може \nt{GLUT}-нейрон видати від'ємний сигнал. Рішення знайоме з розділу~\ref{sec:glutcomputer}: \emph{два проводи}. Кожен суглоб обслуговує пара м'язів, що тягнуть у протилежні боки (рис.~\ref{fig:antagonists}). Якщо їхні сили $F_1$ і $F_2$, а плече $\ell$, то обертальний момент і жорсткість суглоба дорівнюють
\begin{empheq}[box=\keybox]{equation}
M \;=\; \ell\,(F_1-F_2), \qquad K \;\approx\; \kappa_m\,(F_1+F_2),
\label{eq:antagonists}
\end{empheq}
де $\kappa_m$ --- коефіцієнт, що пов'язує силу м'яза з його жорсткістю: напружений м'яз пружинить сильніше за розслаблений. Різниця сил задає момент, сума --- жорсткість~\cite{Hogan1984}. Двома проводами мозок керує двома незалежними величинами: куди повернути суглоб і наскільки твердо його тримати. Коли треба втримати чашку, не розхлюпавши, ми напружуємо обидва м'язи одразу: момент той самий, жорсткість вища, і випадковий поштовх менше збиває руку.

\begin{figure}[t]
\centering
\begin{tikzpicture}[>=Stealth,
  nrn/.style={circle,draw,very thick,minimum size=0.95cm,inner sep=0pt,font=\scriptsize},
  inh/.style={-{Bar[width=7pt]},thick,red!70!black}]
% the joint: a bar on a pivot, two springs pulling down on either side
\fill[gray!30] (4.6,-0.25) -- (5.4,-0.25) -- (5,0.15) -- cycle;
\draw[line width=3pt,gray!50!black] (2.4,0.2) -- (7.6,0.2);
\fill[black] (5,0.2) circle (0.08);
\draw[decorate,decoration={coil,segment length=5pt,amplitude=4pt},very thick,blue!60!black] (3.0,0.2) -- (3.0,-1.6);
\draw[decorate,decoration={coil,segment length=5pt,amplitude=4pt},very thick,orange!85!black] (7.0,0.2) -- (7.0,-1.6);
\draw[very thick] (2.5,-1.6) -- (3.5,-1.6); \draw[very thick] (6.5,-1.6) -- (7.5,-1.6);
\node[font=\small,blue!60!black] at (2.35,-0.75) {$F_1$};
\node[font=\small,orange!85!black] at (7.65,-0.75) {$F_2$};
\draw[->,thick] (5.9,0.75) arc (60:120:1.8) node[above,font=\scriptsize,pos=0.5] {$M=\ell(F_1-F_2)$};
% motor neurons and reflex circuit
\node[nrn,fill=green!12] (M1) at (1.6,-3.0) {\nt{ACET}};
\node[nrn,fill=green!12] (M2) at (8.4,-3.0) {\nt{ACET}};
\node[nrn,fill=red!10] (G) at (5.0,-3.0) {\nt{GLYC}};
\node[nrn,fill=blue!6] (S) at (1.6,-5.0) {\nt{GLUT}};
\draw[->,very thick] (M1) -- (3.0,-1.7);
\draw[->,very thick] (M2) -- (7.0,-1.7);
\draw[->,thick] (S) -- (M1);
\draw[->,thick] (S) -- (G);
\draw[inh] (G) -- (M2);
\draw[->,thick,densely dashed,gray!60!black] (3.15,-1.2) to[bend left=25] node[pos=0.35,right,font=\scriptsize,black] {розтяг} (S.north east);
\draw[->,thick,blue!60!black] (-0.3,-3.0) node[left,font=\scriptsize,black,align=right] {команда\\мозку} -- (M1);
\draw[->,thick,blue!60!black] (10.3,-3.0) node[right,font=\scriptsize,black,align=left] {команда\\мозку} -- (M2);
\end{tikzpicture}
\caption{Два м'язи на один суглоб і найшвидша петля зворотного зв'язку. \nt{ACET}-нейрони отримують команди мозку й тягнуть суглоб у різні боки; різниця сил повертає його, сума робить жорсткішим. Датчик розтягу (\nt{GLUT}) лівого м'яза збуджує його власний \nt{ACET}-нейрон і через \nt{GLYC}-посередника гальмує нейрон м'яза-антагоніста --- заборона з розділу~\ref{sec:gaba}.}
\label{fig:antagonists}
\end{figure}

\section{Зворотний зв'язок із затримкою}

М'язи не працюють наосліп. У кожному сидять датчики розтягу, про які йшлося в таблиці~\ref{tab:sensors}, і найкоротша петля замикається просто в спинному мозку (рис.~\ref{fig:antagonists}). М'яз несподівано розтягнули --- датчик розтягу збуджує його власний \nt{ACET}-нейрон, м'яз скорочується й повертає суглоб. Водночас через гальмівний нейрон-посередник вимикається м'яз-антагоніст, щоб не заважав. Цим посередником працює \nt{GLYC} --- той самий тип із таблиці~\ref{tab:types}, якому не дістався власний розділ: його місце саме тут, у швидких петлях спинного мозку.

Кожна петля має затримку $d$: спинномозкова --- близько $25$--$30\ms$ для руки, петля через кору --- у півтора--три рази більшу, петля через око --- $100$--$200\ms$. А затриманий зворотний зв'язок, як ми бачили в твердженні~\ref{prop:ei}, розгойдує систему. Для найпростішого регулятора, що виправляє відхилення $x$ зі швидкістю $G$ за відомостями $d$-секундної давності, $\dd x/\dd t=-G\,x(t-d)$, умова стійкості така~\cite{Hayes1950}:
\begin{empheq}[box=\keybox]{equation}
G\,d \;<\; \frac{\pi}{2} ,
\label{eq:delaystable}
\end{empheq}
а на межі система коливається з частотою $1/(4d)$. Ми перевірили це на моделі: при $d=30\ms$ відхилення загасає при $G=40$ за секунду, а при $G=55$ за секунду, трохи вище за межу $52$, розгойдується.

Звідси два наслідки. Перший: що довша петля, то повільніше вона мусить виправляти. Через око із затримкою в півтори сотні мілісекунд руку не можна поправляти швидше, ніж приблизно за десяту частку секунди, інакше вона затремтить. Другий: межа стійкості спинномозкової петлі, $1/(4\cdot30\ms)\approx8$ коливань за секунду, близька до частоти дрібного тремтіння, яке є в кожної здорової людини, --- вісім-дванадцять коливань за секунду. Петля розтягу --- одне з імовірних джерел цього тремтіння~\cite{McAuleyMarsden2000}.

Як тоді мозок рухається швидко й точно, якщо зворотний зв'язок запізнюється? За поширеною гіпотезою --- \emph{передбачає}: тримає внутрішню модель тіла й обчислює, яким буде результат команди, не чекаючи на датчики~\cite{WolpertGhahramani2000}. Датчики потрібні, щоб виправляти модель, а не кожен рух. Інженер назвав би це керуванням із випередженням і спостерігачем стану.

\section{Генератори ритму рухів}

Ходьба, дихання, жування --- ритмічні рухи. Чи потрібен для них тактовий генератор? Ні. У спинному мозку й стовбурі є невеликі мережі, які самі породжують ритм, навіть якщо на них не надходить нічого ритмічного~\cite{MarderBucher2001}. Найпростішу таку мережу зібрано з деталей, які в нас уже є: дві групи \nt{GLUT}-нейронів, що гальмують одна одну через \nt{GABA}- або \nt{GLYC}-посередників, як на рис.~\ref{fig:wta}, і повільно втомлюються, як пам'ять розділу~\ref{sec:glutcomputer}:
\begin{empheq}[box=\keybox]{equation}
\tau\,\frac{\dd r_i}{\dd t} = -r_i + \bigl[I - \beta\,r_j - a_i\bigr]_+ ,
\qquad
\tau_a\,\frac{\dd a_i}{\dd t} = -a_i + b\,r_i ,
\label{eq:halfcentre}
\end{empheq}
де $a_i$ --- втома групи $i$, $j$ --- інша група. Взаємне гальмування з $\beta>1$ вибирає переможця (твердження~\ref{prop:wta}). Переможець працює і втомлюється; коли втома з'їдає його вхід, переможений, уже відпочилий, виривається вперед і сам починає втомлюватися. Групи працюють по черзі, і кожна керує своїм м'язом пари (рис.~\ref{fig:halfcentre}).

\begin{figure}[t]
\centering
\begin{tikzpicture}[>=Stealth,x=3.2cm,y=2.4cm]
\draw[->,gray!70!black] (0,0) -- (4.2,0) node[right,font=\small,black] {$t$, с};
\draw[->,gray!70!black] (0,0) -- (0,1.2) node[left,font=\small,black] {$r$};
\foreach \t in {0,1,2,3,4}{\draw[gray!60!black] (\t,-0.03) -- (\t,0.03); \node[below,font=\scriptsize] at (\t,0) {\t};}
\draw[thick,blue!60!black] plot file {figures/halfcentre1.dat};
\draw[thick,orange!85!black] plot file {figures/halfcentre2.dat};
\node[font=\scriptsize,blue!60!black,anchor=west] at (1.6,1.28) {група 1: «вперед»};
\node[font=\scriptsize,orange!85!black,anchor=west] at (2.7,1.28) {група 2: «назад»};
\end{tikzpicture}
\caption{Генератор ритму за моделлю~\eqref{eq:halfcentre}: $I=1$, $\beta=2$, $b=2$, $\tau=20\ms$, $\tau_a=0.4\,\text{с}$. Дві групи із взаємним гальмуванням і втомою працюють по черзі з періодом $0.84\,\text{с}$, хоча на вхід подано сталий сигнал.}
\label{fig:halfcentre}
\end{figure}

У нашій моделі за сталого входу групи змінюють одна одну з періодом $0.84\,\text{с}$, приблизно вдвічі більшим за час утоми $\tau_a$. Стала команда згори --- «іди» --- перетворюється на ритм на місці. Це ще один місцевий ритм, як ритм петлі \nt{GLUT}--\nt{GABA} у розділі~\ref{sec:gaba}: він виникає, лише коли мережа працює, і задає такт тільки тим м'язам, якими керує.

\begin{keyblock}
\begin{proposition}[Вихід машини]
\label{prop:muscles}
Мозок керує м'язами як ієрархія регуляторів. Нагорі вибирається дія (розділ~\ref{sec:dofa}); нижче місцеві генератори перетворюють сталі команди на ритм, а швидкі петлі спинного мозку тримають суглоби. Сила задається з рухомою комою: порядок --- кількістю одиниць, що вмикаються за розміром, мантиса --- частотою імпульсів. Кожним суглобом керує пара тягнучих м'язів, різниця сил яких задає момент, а сума --- жорсткість. Ці пристрої виміряно; що мозок спирається на внутрішню модель тіла, --- добре обґрунтована гіпотеза.
\end{proposition}
\end{keyblock}

\section{Підсумок}

\begin{table}[t]
\centering
\footnotesize
\setlength{\tabcolsep}{4pt}
\renewcommand{\arraystretch}{1.15}
\begin{tabular}{>{\raggedright}p{3.0cm}>{\raggedright}p{4.6cm}>{\raggedright\arraybackslash}p{5.2cm}}
\toprule
Що робить вихід & Як & Що відомо \\
\midrule
перетворює імпульси на силу & сума посмикувань~\eqref{eq:forcesum}, фільтр нижніх частот & виміряно; лінійна сума --- спрощення \\
дозує силу & увімкнення одиниць за розміром, рухома кома~\eqref{eq:sizeprinciple} & порядок увімкнення виміряно \\
штовхає, уміючи лише тягнути & пара м'язів: момент --- різниця, жорсткість --- сума~\eqref{eq:antagonists} & виміряно \\
тримає суглоб & петля розтягу, заборона антагоніста через \nt{GLYC} & виміряно \\
не тремтить & $Gd<\pi/2$~\eqref{eq:delaystable} & випливає з теорії регулювання; внесок у тремтіння --- імовірна причина \\
рухається швидко & передбачення за внутрішньою моделлю & добре обґрунтована гіпотеза \\
ходить і дихає ритмічно & генератор ритму~\eqref{eq:halfcentre} & генератори виміряно; найпростіша схема --- модель \\
\bottomrule
\end{tabular}
\caption{Як машина керує м'язами.}
\label{tab:muscles}
\end{table}

Вихід машини влаштовано тими самими прийомами, що й усе інше (таблиця~\ref{tab:muscles}): два проводи замість знака, фільтр нижніх частот замість ЦАП, логарифмічна шкала замість лінійної, взаємне гальмування й утома замість тактового генератора. Вхід (розділ~\ref{sec:sensors}) і вихід міряють світ у відносних частках, а між ними працює машина, якій присвячено цю книгу.

\section{Задачі}

\begin{exercise}[\lvlA]
\label{ex:13:1}
\emph{Рухома кома в цифрах.} У м'язі $N=100$ рухових одиниць із силами $f_n=f_1q^{\,n-1}$; найсильніша в $100$ разів сильніша за найслабшу. (а)~Знайдіть $q$, відносний крок~\eqref{eq:sizeprinciple} і динамічний діапазон у децибелах. (б)~Який крок за сили $10f_1$ у «лінійного ЦАП» зі $100$ однакових одиниць тієї самої сумарної сили?
\end{exercise}

\begin{exercise}[\lvlA]
\label{ex:13:2}
\emph{Ціна запасу надійності.} На кожен імпульс синапс на м'язовому волокні вивільняє пуассонівську кількість порцій із середнім $m$, а волокно спрацьовує за $n_{\text{пор}}=20$ порцій і більше; запас надійності $S=m/n_{\text{пор}}$. Скільки збоїв на добу дає одиниця, що працює з частотою $20$ імпульсів за секунду, при $S=2$ і при $S=4$?
\end{exercise}

\begin{exercise}[\lvlB]
\label{ex:13:3}
\emph{М'яз як фільтр.} Посмикування~\eqref{eq:twitch} з $T_c=50\ms$ --- імпульсна характеристика лінійної системи. (а)~Знайдіть її передавальну функцію $H(s)$ і частоту зрізу. (б)~За якої частоти регулярних імпульсів розмах пульсацій сили становить $10\%$ від середньої сили?
\end{exercise}

\begin{exercise}[\lvlB]
\label{ex:13:4}
\emph{Швидше за затримку не виправиш.} Регулятор $\dd x/\dd t=-G\,x(t-d)$. (а)~Покажіть, що найшвидше відхилення загасає без коливань при $Gd=1/e$, зі швидкістю $1/d$. (б)~Знайдіть це $G$ і сталу часу загасання для спинномозкової петлі, $d=30\ms$, і для петлі через око, $d=150\ms$.
\end{exercise}

\begin{exercise}[\lvlB]
\label{ex:13:5}
\emph{Період генератора ритму.} При $\tau\ll\tau_a$ період $T$ моделі~\eqref{eq:halfcentre} задається рівнянням $(\beta-1)(1-E^{2+b})/(\beta-E)=b\,(1-E^{1+b})/(1+b)$, де $E=e^{-T/(2\tau_a)}$. (а)~Знайдіть $T$ при $\beta=b=2$, $\tau_a=0.4$~с. (б)~Покажіть, що $T$ не залежить від входу $I$, а ритм можливий лише при $b>\beta-1$.
\end{exercise}

\begin{exercise}[\lvlB]
\label{ex:13:6}
\emph{Моделювання генератора.} Імпортуйте функцію \texttt{halfcentre} з \texttt{models/\allowbreak muscle\_models.py} (сам файл не запускайте) і виміряйте період, відкинувши перші два цикли, при $b=0.8$, $1.05$, $1.2$, $1.5$, $2$, $4$ і при $I=0.5$, $1$, $2$. Чи може команда «іди швидше» змінювати темп через $I$?
\end{exercise}

\begin{exercise}[\lvlB]
\label{ex:13:7}
\emph{Жорсткість проти рефлексу.} Суглоб охоплено петлею розтягу: $\dd\theta/\dd t=-a\theta(t)-G\theta(t-d)$, $a=K/B$, $d=30\ms$. Зведіть її до задачі~\ref{ex:13:4} і знайдіть найменше $a$, за якого відхилення загасає без коливань зі сталою часу $15\ms$, і потрібне $G$. Як змінювати $G$ зі збільшенням спільного напруження м'язів?
\end{exercise}

\begin{exercise}[\lvlC]
\label{ex:13:8}
\emph{Ручка темпу.} За задачами~\ref{ex:13:5} і~\ref{ex:13:6} вхід $I$ генератора~\eqref{eq:halfcentre} змінює лише розмах, але не темп. Запропонуйте мінімальну зміну моделі з деталей книги, за якої нижче від деякого $I_{\min}$ ритму немає, а вище період спадає зі зростанням $I$, хоча б удвічі при потроєнні $I$. Знайдіть $I_{\min}$ при $\tau\ll\tau_a$ і перевірте моделюванням.
\end{exercise}
